% ============================================================
% Notes: Laplace Transform (A&W style)
% Topics: definition, derivatives, delta/Heaviside, convolution,
%         inversion (Bromwich integral)
% Source: Arfken & Weber, Mathematical Methods for Physicists
% ============================================================

\section{Laplace transforms}

\subsection{Definition}
For a function $f(t)$ (typically $t\ge 0$), the Laplace transform is
\begin{equation}
F(s) = \mathcal{L}\{f(t)\}(s) := \int_{0}^{\infty} e^{-st} f(t)\,dt,
\qquad s\in\mathbb{C}.
\end{equation}
The integral converges in a half-plane
\[
\Re(s) > s_0
\]
determined by the growth of $f(t)$ (``region of convergence'').

\paragraph{Linearity.}
\begin{equation}
\mathcal{L}\{af+bg\}=a\mathcal{L}\{f\}+b\mathcal{L}\{g\}.
\end{equation}

\subsection{Useful basic transforms}
\begin{align}
\mathcal{L}\{1\} &= \frac{1}{s},\\
\mathcal{L}\{t^n\} &= \frac{n!}{s^{n+1}}, \qquad n=0,1,2,\dots\\
\mathcal{L}\{e^{at}\} &= \frac{1}{s-a},\\
\mathcal{L}\{\cos(\omega t)\} &= \frac{s}{s^2+\omega^2},\\
\mathcal{L}\{\sin(\omega t)\} &= \frac{\omega}{s^2+\omega^2}.
\end{align}

% ------------------------------------------------------------

\section{Laplace transform of derivatives}

Assume $f(t)$ is sufficiently smooth and of exponential order so that the transforms exist.

\subsection{First derivative}
\begin{align}
\mathcal{L}\{f'(t)\}(s)
&= \int_0^\infty e^{-st} f'(t)\,dt\\
&= \Big[e^{-st}f(t)\Big]_{0}^{\infty}
+ s\int_0^\infty e^{-st}f(t)\,dt\\
&= -f(0) + sF(s).
\end{align}
Hence the key identity:
\begin{equation}
\boxed{\mathcal{L}\{f'(t)\}(s)=sF(s)-f(0).}
\end{equation}

\subsection{Second derivative}
Apply the first-derivative rule to $f'(t)$:
\begin{equation}
\boxed{\mathcal{L}\{f''(t)\}(s)=s^2F(s)-sf(0)-f'(0).}
\end{equation}

\subsection{$n$th derivative}
By induction:
\begin{equation}
\boxed{
\mathcal{L}\{f^{(n)}(t)\}(s)
=
s^nF(s)-s^{n-1}f(0)-s^{n-2}f'(0)-\cdots-f^{(n-1)}(0).
}
\end{equation}

\paragraph{High-yield use.}
Laplace transforms turn ODE initial-value problems into algebra:
\[
\text{ODE in }t \quad \longrightarrow \quad \text{algebraic equation in }s.
\]

% ------------------------------------------------------------

\section{Laplace transform of Dirac delta and Heaviside functions}

\subsection{Dirac delta $\delta(t-a)$}
Using the defining property of the delta distribution,
\[
\int_0^\infty \delta(t-a)\,g(t)\,dt = g(a), \qquad a>0,
\]
we obtain
\begin{equation}
\boxed{
\mathcal{L}\{\delta(t-a)\}(s)
=
\int_0^\infty e^{-st}\delta(t-a)\,dt
=
e^{-as}.
}
\end{equation}
Special case:
\begin{equation}
\mathcal{L}\{\delta(t)\}(s)=1.
\end{equation}

\subsection{Heaviside step function $H(t-a)$}
Define
\[
H(t-a)=
\begin{cases}
0, & t<a,\\
1, & t>a.
\end{cases}
\]
Then
\begin{align}
\mathcal{L}\{H(t-a)\}(s)
&= \int_0^\infty e^{-st}H(t-a)\,dt
= \int_a^\infty e^{-st}\,dt\\
&= \boxed{\frac{e^{-as}}{s}}.
\end{align}

\subsection{Time-shift theorem (with Heaviside)}
If $g(t)=H(t-a)\,f(t-a)$, then
\begin{equation}
\boxed{
\mathcal{L}\{H(t-a)f(t-a)\}(s)=e^{-as}F(s).
}
\end{equation}
This is the standard tool for piecewise forcing functions.

% ------------------------------------------------------------

\section{Convolution theorem (Laplace)}

\subsection{Definition of convolution on $[0,\infty)$}
For functions $f,g$ on $t\ge 0$, define
\begin{equation}
(f*g)(t) := \int_0^t f(\tau)\,g(t-\tau)\,d\tau.
\end{equation}

\subsection{Laplace transform of convolution}
\begin{equation}
\boxed{
\mathcal{L}\{(f*g)(t)\}(s)=F(s)\,G(s).
}
\end{equation}

\subsection{Inverse statement}
\begin{equation}
\boxed{
\mathcal{L}^{-1}\{F(s)G(s)\}(t)=(f*g)(t).
}
\end{equation}

\paragraph{High-yield use.}
Convolution appears naturally when solving linear ODEs/PDEs with forcing
and in Green's function methods.

% ------------------------------------------------------------

\section{Inverse Laplace transform and Bromwich integral}

\subsection{Inverse Laplace transform (Bromwich inversion)}
The inverse Laplace transform is given by the Bromwich integral:
\begin{equation}
\boxed{
f(t) = \mathcal{L}^{-1}\{F(s)\}(t)
=
\frac{1}{2\pi i}\int_{\gamma-i\infty}^{\gamma+i\infty}
e^{st}F(s)\,ds.
}
\end{equation}
Here $\gamma$ is chosen so that the vertical line $\Re(s)=\gamma$
lies in the region of convergence and to the right of all singularities of $F(s)$.

\subsection{Contour interpretation (complex analysis link)}
The Bromwich integral is evaluated using residue calculus by closing the contour
(usually to the left half-plane for $t>0$) and summing residues of $e^{st}F(s)$:
\begin{equation}
f(t) = \sum_{\text{poles }s_k}\operatorname{Res}\left(e^{st}F(s),\,s=s_k\right),
\qquad (t>0),
\end{equation}
provided the closing arc contribution vanishes.

\paragraph{Key point.}
Laplace inversion is essentially a complex Fourier inversion with $s=\sigma+i\omega$,
and the Bromwich contour integral is the rigorous inversion formula.

% ------------------------------------------------------------

\section{High-yield summary}

\begin{itemize}
\item Definition:
\[
F(s)=\mathcal{L}\{f(t)\}=\int_0^\infty e^{-st}f(t)\,dt.
\]
\item Derivatives:
\[
\mathcal{L}\{f'\}=sF-f(0),\qquad
\mathcal{L}\{f''\}=s^2F-sf(0)-f'(0),
\]
\[
\mathcal{L}\{f^{(n)}\}
=
s^nF-\sum_{k=0}^{n-1}s^{n-1-k}f^{(k)}(0).
\]
\item Delta and Heaviside:
\[
\mathcal{L}\{\delta(t-a)\}=e^{-as},\qquad
\mathcal{L}\{H(t-a)\}=\frac{e^{-as}}{s}.
\]
\item Convolution:
\[
(f*g)(t)=\int_0^t f(\tau)g(t-\tau)\,d\tau,
\qquad
\mathcal{L}\{f*g\}=FG.
\]
\item Bromwich inversion:
\[
f(t)=\frac{1}{2\pi i}\int_{\gamma-i\infty}^{\gamma+i\infty}e^{st}F(s)\,ds.
\]
\end{itemize}
